\documentclass[aspectratio=169,11pt]{beamer}

\usetheme{Madrid}
\usecolortheme{default}
\usefonttheme{professionalfonts}
\setbeamertemplate{navigation symbols}{}
\setbeamertemplate{footline}[frame number]

\usepackage{amsmath, amssymb, booktabs, array}

\title{Frontier Methods, New Data, And Evolving Measurement}
\subtitle{Gender Study --- Week 9}
\author{Teaching deck draft}
\date{}

\begin{document}

\begin{frame}
  \titlepage
\end{frame}

\begin{frame}{Central question}
\begin{itemize}
    \item How are new data, designs, and measurement choices changing gender-and-labor research?
    \item Where are the main empirical shortfalls in the existing literature?
    \item How do we turn those shortfalls into research designs?
\end{itemize}
\vspace{0.25cm}
\[
\text{labor object} \rightarrow \text{data family} \rightarrow \text{observed margin}
\rightarrow \text{identification} \rightarrow \text{welfare interpretation}.
\]
\end{frame}

\begin{frame}{Why methods matter for gender and labor}
\begin{itemize}
    \item Gender gaps are not one outcome: employment, hours, applications, wages, promotion, authority, safety, climate, identity, and welfare all matter.
    \item Better data reveal new margins but do not automatically identify mechanisms.
    \item The same substantive question may require different data: admin panels, matched worker-firm records, audits, postings, platforms, surveys, or time use.
    \item Week 9 is a research-career bridge: find the observed margin, then name what remains missing.
\end{itemize}
\end{frame}

\begin{frame}{Bridge from Weeks 2--8}
\scriptsize
\begin{tabular}{p{0.18\linewidth} p{0.32\linewidth} p{0.38\linewidth}}
\toprule
Week & Substantive object & Week 9 methods question \\
\midrule
2 & care, fertility, household allocation & Which event, time-use, or policy variation isolates the care margin? \\
3 & education, aspirations, sorting & Do we observe expectations, applications, or only realized choices? \\
4 & firms, promotion, pay-setting & Is the gap worker sorting or firm treatment? \\
5 & norms and bargaining & Is the norm measured or inferred from outcomes? \\
6 & policy incidence & Which workers, firms, and households bear the incidence? \\
7 & safety and hidden harms & Do we observe exposure, reporting, or labor response? \\
8 & comparative mechanisms & What transports across institutional settings? \\
\bottomrule
\end{tabular}
\end{frame}

\begin{frame}{Descriptive gap vs causal identification}
\begin{columns}[T,onlytextwidth]
\begin{column}{0.47\linewidth}
\textbf{Descriptive gap}
\[
\Delta^Y = E[Y_i\mid G_i=1] - E[Y_i\mid G_i=0]
\]
\begin{itemize}
    \item Locates where gender differences appear.
    \item Requires clarity about outcome, population, and gender measure.
\end{itemize}
\end{column}
\begin{column}{0.47\linewidth}
\textbf{Causal object}
\[
\tau = E[Y_i(1)-Y_i(0)\mid i\in \mathcal{P}]
\]
\begin{itemize}
    \item Requires a counterfactual.
    \item Requires identifying variation and a target population.
\end{itemize}
\end{column}
\end{columns}
\end{frame}

\begin{frame}{Data-family taxonomy}
\scriptsize
\begin{tabular}{p{0.26\linewidth} p{0.31\linewidth} p{0.31\linewidth}}
\toprule
Data family & Margin seen well & Persistent blind spot \\
\midrule
Matched worker--firm & firm premiums, moves, within/between gaps & nonapplications, hidden harms \\
Administrative panels & lifecycle dynamics, policy timing & care burden, climate, beliefs \\
Audits/field experiments & screening, callbacks, evaluation & long-run careers, equilibrium response \\
Surveys/expectations & beliefs, norms, anticipated discrimination & persistence, response selection \\
Time-use data & unpaid work, care, scheduling & firm treatment if unlinked \\
Postings/text & employer demand, language, stated constraints & realized treatment after hire \\
Platform data & high-frequency supply, tasks, pay & outside options, nonplatform careers \\
Identity/climate & lived treatment, belonging, exposure & small cells, privacy, comparability \\
\bottomrule
\end{tabular}
\end{frame}

\begin{frame}{Setting-to-methods map}
\scriptsize
\begin{tabular}{p{0.25\linewidth} p{0.29\linewidth} p{0.34\linewidth}}
\toprule
Empirical setting & Observed margin & Practical tools \\
\midrule
Childbirth or family shock & earnings, hours, employment event time & event studies, dynamic treatment effects \\
Matched worker--firm panel & worker moves, wages, firm identifiers & worker and firm fixed effects, mover designs \\
Resume or hiring audit & callback, interview, offer & randomized signals, heterogeneity \\
Job-board reform & applications, callbacks, match composition & difference-in-differences, ad-level event studies \\
Time-use diary & care time, market hours & time-budget decompositions \\
Platform panel & hours, timing, tasks, earnings & within-worker decompositions \\
Climate or harassment linkage & complaints, exits, wages, turnover & linked survey-admin, hazards, selection diagnosis \\
\bottomrule
\end{tabular}
\end{frame}

\begin{frame}{Matched employer--employee methods}
\begin{itemize}
    \item Core question: how much of the gender wage gap reflects firm sorting versus firm treatment?
    \item Observed margin: worker-firm matches, wages, mobility, tenure, sometimes occupation or establishment.
    \item Identifying variation: movers across firms and changes in pay around firm transitions.
    \item Tools: AKM-style worker and firm fixed effects, mover designs, within-firm and between-firm decompositions.
\end{itemize}
\vspace{0.2cm}
\[
\bar w_m-\bar w_f =
\sum_j (s_{mj}-s_{fj})\bar w_j
+ \sum_j s_{fj}(\bar w_{mj}-\bar w_{fj}).
\]
\end{frame}

\begin{frame}{Administrative event-study methods}
\begin{itemize}
    \item Core question: how do earnings, hours, employment, and firm outcomes evolve around childbirth, leave, policy, or job transitions?
    \item Observed margin: event time in administrative panels.
    \item Identifying variation: timing around an event or reform.
    \item Tools: event studies, stacked event studies, dynamic treatment effects, outcome decompositions.
    \item Threats: anticipation, endogenous timing, selected event populations, and hidden care or climate margins.
\end{itemize}
\end{frame}

\begin{frame}{Audits and field experiments}
\begin{itemize}
    \item Core question: does the screening or evaluation stage treat otherwise comparable applicants differently?
    \item Observed margin: callback, interview, offer, evaluation score, application response.
    \item Identifying variation: randomized names, gender signals, identity signals, information, or recruiting messages.
    \item Strength: clean treatment variation at a specific stage.
    \item Limit: may not observe later retention, promotion, workplace treatment, or equilibrium adjustment.
\end{itemize}
\end{frame}

\begin{frame}{Applications, postings, and text}
\begin{itemize}
    \item Application data reveal the pre-match margin: who applies to which jobs before accepted jobs are observed.
    \item Posting data reveal employer demand: stated restrictions, flexibility, pay, tasks, seniority, culture signals, and language.
    \item Text-as-data tools include dictionaries, supervised classifiers, validated language measures, and ad-level fixed effects.
    \item Identifying variation can come from platform rule changes, posting bans, randomized ad text, or quasi-experimental ad exposure.
\end{itemize}
\end{frame}

\begin{frame}{Platform data and hidden margins}
\begin{itemize}
    \item Platform records observe high-frequency behavior: hours, timing, location, task choice, speed, ratings, and earnings.
    \item Within-worker decompositions can separate hours, experience, task selection, and scheduling.
    \item The design often sees behavior conditional on platform participation.
    \item Hidden margins remain: outside options, safety, harassment, care constraints, algorithmic exposure, and nonplatform career effects.
\end{itemize}
\end{frame}

\begin{frame}{Incidence vs hidden harms}
\begin{columns}[T,onlytextwidth]
\begin{column}{0.47\linewidth}
\textbf{Incidence}
\begin{itemize}
    \item Who bears the effect?
    \item Workers, firms, coworkers, applicants, households, consumers.
    \item Example margins: wages, hiring, compliance, sorting, hours.
\end{itemize}
\end{column}
\begin{column}{0.47\linewidth}
\textbf{Hidden harms}
\begin{itemize}
    \item What does the direct outcome miss?
    \item Fear, harassment, retaliation, concealment, mobility avoidance, dignity.
    \item A wage effect is not a complete welfare object.
\end{itemize}
\end{column}
\end{columns}
\end{frame}

\begin{frame}{Measurement of gender and identity}
\scriptsize
\begin{tabular}{p{0.26\linewidth} p{0.32\linewidth} p{0.30\linewidth}}
\toprule
Measurement object & Why it matters & Main challenge \\
\midrule
Legal sex/admin category & linkage and long panels & may not match identity or treatment \\
Self-identified gender & lived identity and experience & small cells, privacy, nonresponse \\
Perceived gender & hiring, promotion, harassment mechanism & rarely observed directly \\
Workplace climate & retention, authority, welfare & selective and imperfect measures \\
Intersectional categories & realistic heterogeneity & precision and disclosure risk \\
Cross-time categories & trend interpretation & forms and categories change \\
\bottomrule
\end{tabular}
\end{frame}

\begin{frame}{Legal categories vs lived treatment}
\begin{itemize}
    \item Administrative data may record legal sex because that is what the form collects.
    \item Employers and coworkers may respond to perceived gender, pregnancy, parenthood, name, presentation, sexuality, race-gender stereotypes, or disclosure.
    \item The data category and the mechanism may differ.
    \item A strong paper states which object is measured and which object is theoretically relevant.
\end{itemize}
\end{frame}

\begin{frame}{Empirical challenges and shortfalls}
\begin{itemize}
    \item Selection into observed employment, firms, applications, surveys, and reports.
    \item Sorting versus treatment in occupations, firms, teams, policies, and platforms.
    \item Anticipation before childbirth, leave, promotion, reporting, and exits.
    \item Underreporting and hidden harms in safety, climate, and harassment research.
    \item Firm adaptation to transparency, quotas, compliance, and anti-discrimination rules.
    \item Small cells, privacy, and category changes in identity and intersectional research.
\end{itemize}
\end{frame}

\begin{frame}{Research openings}
\scriptsize
\begin{tabular}{p{0.29\linewidth} p{0.56\linewidth}}
\toprule
Opening & Publishable contribution \\
\midrule
Hidden harms & connect safety, climate, or dignity to labor behavior and welfare \\
Firm-side measurement & observe tasks, supervisors, evaluations, internal ladders \\
Identity/treatment mismatch & distinguish legal category, self-identity, perception, and treatment \\
Search and applications & observe nonapplication, callback, match, and sorting margins \\
Comparative portability & identify what mechanism travels across institutions \\
Intersectional gender & use privacy-safe aggregation with honest precision limits \\
Policy as mechanism test & decompose which margin moves and why \\
\bottomrule
\end{tabular}
\end{frame}

\begin{frame}{Week 9 lab}
\begin{itemize}
    \item Reproduce: synthetic application-gap decompositions inspired by Fluchtmann, Glenny, Harmon, and Maibom.
    \item Diagnose: what application data see and what they miss; compare to matched employer--employee and job-posting designs.
    \item Transfer: platform within-worker decompositions and posting-policy design logic.
    \item Optional memo: how evolving identity measurement changes audit or application designs.
\end{itemize}
\end{frame}

\begin{frame}{Bridge to Week 10}
\begin{enumerate}
    \item State the labor object.
    \item Define the gender concept.
    \item Name the observed margin.
    \item Identify the variation.
    \item Diagnose sorting, treatment, and hidden harms.
    \item State the welfare object and limits.
    \item Say what readers should update.
\end{enumerate}
\end{frame}

\begin{frame}{Takeaways}
\begin{itemize}
    \item New data expand the gender-and-labor frontier by making margins visible.
    \item The research design still turns on identification, measurement, and welfare interpretation.
    \item The best projects convert empirical shortfalls into disciplined contributions.
    \item Week 10 uses this toolkit to build student research designs.
\end{itemize}
\end{frame}

\end{document}
